%! AP Statistics — Unit 2, Lesson 2.8 — Exit Ticket — Answer Key
\documentclass[10pt]{article}
\usepackage{apstats-article}
\usepackage{apstats-key}

\begin{document}

\pageheader{Unit 2, Lesson 2.8}{Exit Ticket --- Answer Key}
\namedateperiod

\vspace{0.18in}

% ── SCENARIO ─────────────────────────────────────────────────────────────────
\begin{tcolorbox}[colback=sky, colframe=navy, boxrule=0.9pt, arc=2mm,
    left=4mm, right=4mm, top=2.5mm, bottom=2.5mm]
\small\textit{A researcher studying coffee shops records average daily customers ($x$) and
daily revenue in hundreds of dollars ($y$) for 35 locations.
Summary statistics: $\bar{x} = 120$ customers, $\bar{y} = 18$ hundred dollars,
$s_x = 40$ customers, $s_y = 5$ hundred dollars, $r = 0.80$.}
\end{tcolorbox}

\vspace{0.16in}

\begin{enumerate}[leftmargin=1.6em, itemsep=0.16in]

  \item Compute the \textbf{slope} $b$ and \textbf{intercept} $a$. Show all work.\\[8pt]
        \ansline{$b = r\,\dfrac{s_y}{s_x} = 0.80\cdot\dfrac{5}{40} = 0.80 \times 0.125 = 0.10$ hundred dollars per customer}\\[3pt]
        \ansline{$a = \bar{y} - b\bar{x} = 18 - (0.10)(120) = 18 - 12 = 6$ hundred dollars}

  \item Write the regression equation using the \textbf{names of the variables}.\\[8pt]
        \ansline{$\widehat{\text{revenue}} = 6 + 0.10(\text{customers})$}

  \item Compute $r^2$ and write the full AP-template \textbf{interpretation} in context.\\[8pt]
        \ansline{$r^2 = (0.80)^2 = 0.64$}\\[3pt]
        \ansline{About 64\% of the variation in daily revenue is explained by the least-squares}\\[3pt]
        \ansline{regression of daily revenue on average daily customers.}

\end{enumerate}

\vspace{0.14in}

% ── AP-STYLE MC ──────────────────────────────────────────────────────────────
\begin{tcolorbox}[colback=redbg, colframe=redacc, boxrule=0.8pt, arc=2mm,
    left=4mm, right=4mm, top=2.5mm, bottom=2.5mm]
\small\textbf{AP-Style Practice}\enspace
Computer output for a regression shows $r^2 = 0.64$. Which of the following is the
\emph{best} interpretation of this value?

\vspace{7pt}
\begin{enumerate}[label=\textbf{(\Alph*)}, leftmargin=2em, itemsep=4pt]
  \item The correlation between the two variables is $0.64$.
  \item 64\% of the data points are close to the regression line.
  \item About 64\% of the variation in the response variable is explained by the least-squares regression on the explanatory variable.
        \textcolor{keyred}{\textbf{$\leftarrow$ correct}}
  \item The regression line passes through 64\% of the data points.
\end{enumerate}
\end{tcolorbox}

\vspace{0.08in}
\noindent\textcolor{slate}{\small My answer: \ans{C} \quad Because: \ans{(C) uses the correct AP template: ``variation in response,'' ``explained by.'' (A) confuses $r^2$ with $r$. (B) and (D) describe proximity, not explained variation.}}

\begin{teachernote}
Items 1--2 are the core computation; all students should reach these. Item 3 (interpretation)
is the AP-exam skill. Collect the exit ticket and use it to sort students: those who write
``the data are 64\% accurate'' need direct coaching on the AP template before Lesson 2.9.
\end{teachernote}

\end{document}
